\section{Linear acquisition on a protected component}
\label{sec:linear-moment-acquisition}
The preceding all-node construction is deterministic in its nominal
mesh and remains a valid acquisition rule. Once a complete expanded
component is located, its prefix can instead be stopped geometrically.
This makes the inverse linear and permits direct estimation of its
integrals. In particular, one need not estimate the occupation at
every point of a fine mesh. Throughout this subsection $a=te_1$ is
the fixed command vector; $\tau$ denotes a moment tolerance.

\begin{lemma}[A geometrically stopped prefix]
\label{lem:protected-linear-prefix}
Let $P=C_0-A_0$ be one complete expanded component, separated from
all other expanded components by at least $g>t$. Suppose a known
compact rational convex polygon $Q$ satisfies
\begin{equation}\label{eq:protected-outer-polygon}
 P\subset Q\subset P+\rho\overline B,
                 \qquad \rho+t<g,
\end{equation}
where $\overline B$ is the closed Euclidean unit disk.
Choose a known integer $K$ with $Kt>\diam Q$. For $x\in Q$ put
\[
 m_Q(x)=\min\{j\ge1:x+ja\notin Q\},
\]
and put $m_Q(x)=0$ for $x\notin Q$. Then $m_Q(x)\le K$, and the
component occupation has the pointwise representation
\begin{equation}\label{eq:protected-linear-prefix}
 v_C(x)=\sum_{k=0}^{K-1}\one_{\{k<m_Q(x)\}}
       \bigl(F_-(x+(k+1)a)-F_+(x+ka)\bigr).
\end{equation}
At rational $x$, the integer $m_Q(x)$ is found by finitely many exact
rational polygon membership tests. It does not depend on observed bits.
\end{lemma}
\begin{proof}
If $x\notin Q$, then $x\notin P$ and both sides vanish. Suppose
$x\in Q$. Since $Kt>\diam Q$, the point $x+Ka$ is outside $Q$,
so $1\le m_Q(x)\le K$. Put $y=x+m_Q(x)a$. By the definition of
first exit, $y-a\in Q$, whence
\[
                 \dist(y,P)\le\rho+t<g.
\]
The point $y$ is not in $P$, since it is outside $Q$; the displayed
strict inequality also excludes every other expanded component.
Thus $v(y)=0$. Furthermore $Q$ meets no other expanded component,
so $v(x)=v_C(x)$. Sum the endpoint increment
\eqref{eq:two-field-occupation-increment} from $x$ to $y$ with its
sign reversed. This gives \eqref{eq:protected-linear-prefix}.
The argument uses the physical mean fields and the closed-solid
convention, including atoms charging boundary starts. For rational
vertices, membership in the closed polygon is a finite family of
rational linear inequalities. Enumeration through $K$ gives the
claimed exact rule, with boundary membership resolved without a
numerical tolerance.
\end{proof}

The polygon required by the lemma follows from the existing geometric
output. Indeed, let $\widehat P$ be a rational convex polygon with
$d_H(\widehat P,P)\le\xi$. Fix a positive dyadic $r$ so small that
$3r+t<g$, and require $\xi\le r$. Then
\begin{equation}\label{eq:protected-polygon-construction}
 Q=\widehat P+[-r,r]^2,\qquad
 P\subset Q\subset P+3r\overline B.
\end{equation}
The first inclusion uses $P\subset\widehat P+\xi\overline B$;
the second uses $\xi+\sqrt2r\le3r$. Thus $\rho=3r$ works.
Only fixed coarse accuracy is needed for this construction. The
prior diameter bound fixes $K$ independently of $\tau$ and of the
realized geometric estimate. Complete-component acquisition and its
error reserve, not a component clipped by an observation window,
are essential here. The component and the fixed prefix extension
must both lie in the protected acquisition aperture; enlarging the
prior-dependent aperture by this fixed collar suffices.

\begin{theorem}[Shared moments from bounded two-bit records]
\label{thm:linear-moment-sampling}
Assume \eqref{eq:protected-outer-polygon}. Fix a dyadic square $W$
containing $Q$ with a collar, and choose the rational normalization
$R$ so that $W\subset(-R/2,R/2)^2$. Let $n\ge1$,
$0<\tau<1/4$, and $0<\delta<1/4$. A randomized finite controller
using only $a$ and $-a$ returns rational estimates of all
\[
 I_\alpha=\int v_C(x)x_R^\alpha\dd x,
                   \qquad |\alpha|\le n,
\]
with simultaneous errors at most $\tau$ and probability at least
$1-\delta$. Conditional on the stated geometry, its attempted-bit
bound is
\begin{equation}\label{eq:linear-moment-cost}
 N_{\rm lin}\le
 2\left\lceil 8|W|^2K^2\tau^{-2}
                 \log\frac{2d_n}{\delta}\right\rceil,
        \qquad d_n=\binom{n+2}{2}.
\end{equation}
All moments use the same records. The nominal mesh size is comparable
to $\tau/(n+1)$, but the controller samples mesh points rather than
visiting the entire mesh. Distinct sites and command occurrences
are at most $N_{\rm lin}$, and their coordinate description length
per occurrence is $C\log(C(n+1)/\tau)$ under the fixed aperture
bound. Controller random-bit generation, polygon arithmetic,
evaluation and storage of moments, physical positioning and travel
are not included in the attempted-bit count. The law may be any
compact probability satisfying the support and separation assumptions;
no density upper bound, lower bound or regularity modulus is required
for this conditional moment stage.
\end{theorem}
\begin{proof}
Partition $W$ into equal dyadic squares of side $\ell$, and let
$G_\ell$ consist of their centers. Take $\ell$ comparable to
$\tau/(n+1)$ with a sufficiently small fixed constant. At each
independent block choose $X$ uniformly from $G_\ell$ and $L$
uniformly from $\{0,\ldots,K-1\}$, independently. Make two fresh
preparations: a bit $Y_+$ at nominal center $X+La$ with command $a$,
and a bit $Y_-$ at nominal center $X+(L+1)a$ with command $-a$.
Both are charged, even when the coefficient in the following record
vanishes. Form
\begin{equation}\label{eq:linear-two-bit-record}
 Z=K\one_{\{L<m_Q(X)\}}(Y_--Y_+),
           \qquad T_\alpha=|W|X_R^\alpha Z.
\end{equation}
Conditioning first on $X,L$ and then averaging over $L$ gives
\begin{equation}\label{eq:linear-record-unbiasedness}
 \E(Z\mid X)=v_C(X),\qquad
 \E T_\alpha=\ell^2\sum_{x\in G_\ell}x_R^\alpha v_C(x).
\end{equation}
The two hidden launches are independent fresh draws from the same
stationary law; equality of their realized displacements is neither
required nor observed. The signed record can be negative. Clipping
it to a probability would destroy the displayed unbiasedness and
is not part of the rule.

We bound the difference between the last Riemann sum and $I_\alpha$.
For each fixed $z\in A_0$, the indicator $\one_{C_0-z}$ has support
in $P\subset\operatorname{int}W$. The total area of mesh cells
meeting its boundary is at most $C\ell$, uniformly in $z$, by the
uniform perimeter and diameter bounds for a convex body. On all
other cells the normalized monomial has magnitude at most one and
Lipschitz constant at most $C(n+1)$. The quadrature error is therefore
at most $C(n+1)\ell$. Boundary membership affects only the boundary
cells. Integration against the arbitrary probability $\mu_0$
preserves this estimate, by the finite sum and Tonelli's theorem.
Our choice of mesh makes the bias at most $\tau/2$. No regularity
of the discontinuous coefficients $m_Q$ is needed: the expectation
in \eqref{eq:linear-record-unbiasedness} is the original occupation.

For every $\alpha$, $|T_\alpha|\le |W|K$. For $q$ independent
blocks, the bounded-variable exponential estimate gives
\[
 \Pr\left(\left|q^{-1}\sum_{j=1}^qT_{\alpha,j}
                  -\E T_\alpha\right|>\tau/2\right)
       \le 2\exp\left(-\frac{q\tau^2}{8|W|^2K^2}\right).
\]
Taking $q$ to be the ceiling in \eqref{eq:linear-moment-cost} and
summing these probabilities over the $d_n$ multi-indices proves the
claim, after adding the quadrature bias. The samples for different
multi-indices need not be independent; the union bound uses only
the independent blocks. Every block has exactly two attempted bits.
The signed rational records and their averages are rational.

Each center is a mesh center plus one of finitely many fixed dyadic
translates $ka$. Thus its description has the asserted length. The
number of available grid points does not multiply the number of
observations: an exact uniform finite choice determines a target,
and no bit at an unchosen target is acquired. This is a finite
controller using only computational randomization. Conditional on
any earlier valid geometric output the same bounds hold, since all
random choices and preparations at this stage are fresh.
\end{proof}

\begin{corollary}[An improved separated budget for the joint inverse]
\label{cor:linear-joint-budget}
Under the hypotheses and notation of
Theorem~\ref{thm:two-field-finite-law}, its geometric, probability,
discrete and prediction conclusions hold with the alternative bound
\begin{equation}\label{eq:linear-joint-budget}
 N\le N_{\rm geom}(c a_m,\delta/2)
                   +Ca_m^{-2}\log\frac{Cm}{\delta}.
\end{equation}
The two command vectors, the positive rational probability output
with at most $\binom{4m+2}{2}$ atoms, and the coordinate-description
bound \eqref{eq:two-field-joint-description} are unchanged. The
mesh of candidate output atoms and deterministic optimization costs
are not observation counts. The exponential sufficient bound
\eqref{eq:two-field-joint-exponential-cost} remains valid; no
parametric rate for the complete joint inverse is asserted.
\end{corollary}
\begin{proof}
Run the geometric stage with failure allowance $\delta/2$ and
accuracy $c a_m$. On its success event choose the complete component
and form \eqref{eq:protected-polygon-construction}, with the fixed
$r$ and a prior-dependent $K$. For small $a_m$ its certified error
is at most $r$. Apply Theorem~\ref{thm:linear-moment-sampling} with
$n=4m$, tolerance $c a_m$ and failure allowance $\delta/2$.
It replaces only the acquisition of the unnormalized moments in
the proof of Theorem~\ref{thm:two-field-finite-law}.

For clarity, the remaining factorization is the positive construction
in Appendix~\ref{sec:moment-factor-details}: use the exact moments
of the actual rational polygon factor, divide the estimated moments
by its positive area, and solve the rational nonnegative fitting
programme on the padded footprint grid. The bounds for factor,
normalization, grid and fitting errors are unchanged. The moment
conditioning and \eqref{eq:two-field-moment-tolerance} yield the
same $W_1$ error. Corollary~\ref{cor:moment-positive-compression}
compresses the positive output without new observations. A union
bound over the geometric event and the conditional moment event gives
probability at least $1-\delta$. Equation~\eqref{eq:linear-moment-cost}
gives \eqref{eq:linear-joint-budget}. Its two terms, including the
geometric cost, are both bounded by the existing exponential bound.
The prediction arguments depend only on the geometric and probability
errors, so their conclusions are unchanged.
\end{proof}

The improvement concerns the separated moment acquisition term. It
does not remove the small moment tolerance, the conditioning of the
positive factorization, or the fine geometry needed to supply that
factor. The all-node estimator above is retained as a deterministic
alternative. Both use exactly the same physical command alphabet
and the same unknown stationary preparation law.
